\documentclass[10pt,a4paper]{amsart}
\usepackage{amsmath,amssymb,amsthm,graphicx}
\usepackage{verbatim}
\newenvironment{reference}{\comment}{\endcomment}
%\newenvironment{reference}{}{}
\newenvironment{slogan}{\comment}{\endcomment}
\newenvironment{history}{\comment}{\endcomment}
% For commutative diagrams we use Xy-pic
\usepackage[all]{xy}

% We use 2cell for 2-commutative diagrams.
\xyoption{2cell}
% Theorem environments.
%
\theoremstyle{plain}
\newtheorem{theorem}[subsection]{Theorem}
\newtheorem{proposition}[subsection]{Proposition}
\newtheorem{lemma}[subsection]{Lemma}

\theoremstyle{definition}
\newtheorem{definition}[subsection]{Definition}
\newtheorem{example}[subsection]{Example}
\newtheorem{exercise}[subsection]{Exercise}
\newtheorem{situation}[subsection]{Situation}

\theoremstyle{remark}
\newtheorem{remark}[subsection]{Remark}
\newtheorem{remarks}[subsection]{Remarks}

\numberwithin{equation}{subsection}

% Macros
%
\def\lim{\mathop{\mathrm{lim}}\nolimits}
\def\colim{\mathop{\mathrm{colim}}\nolimits}
\def\Spec{\mathop{\mathrm{Spec}}}
\def\Hom{\mathop{\mathrm{Hom}}\nolimits}
\def\Ext{\mathop{\mathrm{Ext}}\nolimits}
\def\SheafHom{\mathop{\mathcal{H}\!\mathit{om}}\nolimits}
\def\SheafExt{\mathop{\mathcal{E}\!\mathit{xt}}\nolimits}
\def\Sch{\mathit{Sch}}
\def\Mor{\mathop{\mathrm{Mor}}\nolimits}
\def\Ob{\mathop{\mathrm{Ob}}\nolimits}
\def\Sh{\mathop{\mathit{Sh}}\nolimits}
\def\NL{\mathop{N\!L}\nolimits}
\def\CH{\mathop{\mathrm{CH}}\nolimits}
\def\proetale{{pro\text{-}\acute{e}tale}}
\def\etale{{\acute{e}tale}}
\def\QCoh{\mathit{QCoh}}
\def\Ker{\mathop{\mathrm{Ker}}}
\def\Im{\mathop{\mathrm{Im}}}
\def\Coker{\mathop{\mathrm{Coker}}}
\def\Coim{\mathop{\mathrm{Coim}}}

% Boxtimes
%
\DeclareMathSymbol{\boxtimes}{\mathbin}{AMSa}{"02}

%
% Macros for moduli stacks/spaces
%
\def\QCohstack{\mathcal{QC}\!\mathit{oh}}
\def\Cohstack{\mathcal{C}\!\mathit{oh}}
\def\Spacesstack{\mathcal{S}\!\mathit{paces}}
\def\Quotfunctor{\mathrm{Quot}}
\def\Hilbfunctor{\mathrm{Hilb}}
\def\Curvesstack{\mathcal{C}\!\mathit{urves}}
\def\Polarizedstack{\mathcal{P}\!\mathit{olarized}}
\def\Complexesstack{\mathcal{C}\!\mathit{omplexes}}
% \Pic is the operator that assigns to X its picard group, usage \Pic(X)
% \Picardstack_{X/B} denotes the Picard stack of X over B
% \Picardfunctor_{X/B} denotes the Picard functor of X over B
\def\Pic{\mathop{\mathrm{Pic}}\nolimits}
\def\Picardstack{\mathcal{P}\!\mathit{ic}}
\def\Picardfunctor{\mathrm{Pic}}
\def\Deformationcategory{\mathcal{D}\!\mathit{ef}}
\usepackage{cleveref}
\crefname{originallemma}{original lemma}{original lemmas}
\Crefname{originallemma}{Original lemma}{Original lemmas}
\makeatletter
\newlabel{lemma-closed-in-quasi-compact}{{lemma-closed-in-quasi-compact}{source}{}{sourcepointer-1}{}}
\newlabel{lemma-closed-in-quasi-compact@cref}{{[originallemma][0][]lemma-closed-in-quasi-compact}{[1][0][]source}{}{}{}}
\makeatother
\crefname{originallemma}{original lemma}{original lemmas}
\Crefname{originallemma}{Original lemma}{Original lemmas}
\makeatletter
\newlabel{lemma-constructible-hausdorff-quasi-compact}{{lemma-constructible-hausdorff-quasi-compact}{source}{}{sourcepointer-2}{}}
\newlabel{lemma-constructible-hausdorff-quasi-compact@cref}{{[originallemma][0][]lemma-constructible-hausdorff-quasi-compact}{[1][0][]source}{}{}{}}
\makeatother
\crefname{originallemma}{original lemma}{original lemmas}
\Crefname{originallemma}{Original lemma}{Original lemmas}
\makeatletter
\newlabel{lemma-generic-point-in-constructible}{{lemma-generic-point-in-constructible}{source}{}{sourcepointer-3}{}}
\newlabel{lemma-generic-point-in-constructible@cref}{{[originallemma][0][]lemma-generic-point-in-constructible}{[1][0][]source}{}{}{}}
\makeatother
\title{A subset of a spectral space closed in the constructible topology is spectral}
\author{The Stacks Project Authors}
\date{stacks-project:topology.tex:local-snapshot-20261012; GNU FDL 1.2 or later}
\begin{document}
\maketitle
\noindent\textit{Editorial note.} All delivered mathematical text is the author's own, copied from the pinned Stacks Project chapter source sources/stacks-topology/pinned-original-source.tex; no statement or proof step was rewritten, reordered or removed.\par\medskip
\noindent\textit{Editorial note.} Preamble provenance: the Stacks Project keeps one single shared preamble for every chapter; this card takes it directly from the staged sources/stacks-topology/preamble.tex as three verbatim spans (lines 8-12, 14-18 and 154-217), so no derived file was written and nothing was added or redefined. The wrapper class and the base packages are supplied by the assembler.\par\medskip
\noindent\textit{Editorial note.} Bibliography provenance: neither the counted statement nor its proof contains a citation (this lane rescans the whole span before writing the card), so the card carries no bibliography and no entry was reconstructed.\par\medskip
\noindent\textit{Editorial note.} Selection provenance: this card was opened under the file-based control contract (CONTROL.md), with the mechanical pre-card gate run on its statement and proof interval before the card was written.\par\medskip
\noindent\textit{Allowed standard prerequisites (editorial).}\begin{itemize}
\item The statement that a constructible subset meeting the closure of an irreducible subset contains an open of it. Invoked inside the counted proof (source lines 4658-4660). Shows the generic point of the closure must lie in the subset.
\item The statements that the constructible topology is Hausdorff and quasi-compact and that a closed subset of a quasi-compact space is quasi-compact. Invoked inside the counted proof (source lines 4666-4669). Gives quasi-compactness of the subset in the constructible topology.
\end{itemize}\medskip
\ifdefined\hypertarget\hypertarget{sourcepointer-1}{}\fi
\noindent\textit{Source reference.} Original lemma lemma-closed-in-quasi-compact: sources/stacks-topology/pinned-original-source.tex, source lines 1879--1879. Navigation pointer only: the counted text refers to the Stacks item lemma-closed-in-quasi-compact, which is an existing labeled result of this same chapter and is not part of this card. The pointer carries the target's own label as its number and its definition line as its span; its content is not reproduced here.\par\medskip
\ifdefined\hypertarget\hypertarget{sourcepointer-2}{}\fi
\noindent\textit{Source reference.} Original lemma lemma-constructible-hausdorff-quasi-compact: sources/stacks-topology/pinned-original-source.tex, source lines 4522--4522. Navigation pointer only: the counted text refers to the Stacks item lemma-constructible-hausdorff-quasi-compact, which is an existing labeled result of this same chapter and is not part of this card. The pointer carries the target's own label as its number and its definition line as its span; its content is not reproduced here.\par\medskip
\ifdefined\hypertarget\hypertarget{sourcepointer-3}{}\fi
\noindent\textit{Source reference.} Original lemma lemma-generic-point-in-constructible: sources/stacks-topology/pinned-original-source.tex, source lines 3048--3048. Navigation pointer only: the counted text refers to the Stacks item lemma-generic-point-in-constructible, which is an existing labeled result of this same chapter and is not part of this card. The pointer carries the target's own label as its number and its definition line as its span; its content is not reproduced here.\par\medskip
\begin{lemma}
\label{lemma-spectral-sub}
Let $X$ be a spectral space. Let $E \subset X$ be closed in the constructible
topology (for example constructible or closed). Then $E$ with the induced
topology is a spectral space.
\end{lemma}

\begin{proof}
Let $Z \subset E$ be a closed irreducible subset. Let $\eta$ be the generic
point of the closure $\overline{Z}$ of $Z$ in $X$. To prove that $E$
is sober, we show that $\eta \in E$. If not, then since $E$ is closed
in the constructible topology, there exists a constructible subset
$F \subset X$ such that $\eta \in F$ and $F \cap E = \emptyset$.
By Lemma \ref{lemma-generic-point-in-constructible} this implies
$F \cap \overline{Z}$ contains a nonempty open subset of $\overline{Z}$.
But this is impossible as $\overline{Z}$ is the closure of $Z$ and
$Z \cap F = \emptyset$.

\medskip\noindent
Since $E$ is closed in the constructible topology, it is quasi-compact
in the constructible topology
(Lemmas \ref{lemma-closed-in-quasi-compact} and
\ref{lemma-constructible-hausdorff-quasi-compact}). Hence a fortiori it is
quasi-compact in the topology coming from $X$. If $U \subset X$
is a quasi-compact open, then $E \cap U$ is closed in the constructible
topology, hence quasi-compact (as seen above). It follows that the
quasi-compact open subsets of $E$ are the intersections $E \cap U$
with $U$ quasi-compact open in $X$. These form a basis for the topology.
Finally, given two $U, U' \subset X$ quasi-compact opens, the intersection
$(E \cap U) \cap (E \cap U') = E \cap (U \cap U')$ and $U \cap U'$
is quasi-compact as $X$ is spectral. This finishes the proof.
\end{proof}

\end{document}
